\section{Introduction}\label{sec:introduction}

Cross-national survey programs commit to their items long before they learn how those items behave. A question about trust must be translated, pretested and fielded in many countries before a researcher sees whether its country differences follow the construct it was written for. Language models offer a cheap first reading of such items \citep{argyle2023,horton2023,buskirk2026}. A reading that flagged in advance which new items separate countries would be useful at exactly this stage. We build one such reading and check, link by link, what it does against a benchmark survey and against the respondents of a second survey.

The usual route to a synthetic reading goes through individuals. Silicon sampling describes a person in the prompt, collects the model's answer, and averages many such persons into a country figure \citep{argyle2023}. That route needs two things that are hard to supply. First, it needs a model of each simulated person. The errors of those simulated people are correlated, because one model with one set of assumptions produces all of them. Averaging more of them therefore need not cancel the bias. Second, it needs a sampling frame: someone must decide which people to simulate and in what proportions. The record so far is cautionary. Synthetic responses are sensitive to prompt wording, compress the variance of human answers and depend on the order and labelling of answer options \citep{bisbee2024,boelaert2025,wang2025,dominguezolmedo2024}. Alignment procedures change how faithfully models simulate attitudes \citep{lyman2025}, and methodological work in this journal concludes that validation against human data remains essential \citep{kozlowski2025,broska2025}. Most empirical work asks whether a model reproduces answers that people already gave, whether it is prompted \citep{santurkar2023,durmus2024} or fine-tuned on survey response distributions \citep{cao2025,suh2025}. In the fine-tuned case the human survey sits inside the training signal. Language-model pretesting has also been studied as flaw detection in question wording \citep{sturgis2026}. Whether a method can say which new items will separate countries in an intended direction is a different test. It needs a decision rule fixed before the human criterion is read.

We take a route that never models individuals. Cross-national research usually wants a country-level quantity, such as how countries rank on generalized trust. A probability survey such as the Global Preferences Survey \citep[GPS;][]{falk2018} has already aggregated people into such quantities, with sampling weights. We take its country-level moments as declared inputs and ask the model only for a group-level reading of new item text. In other words, the survey aggregates people, and the language model reads items at the level of the group. An order of countries can survive this design without a representative respondent. It requires only that the instrument's score rise with the group's latent preference. The design has a price. A single choice policy per country has no spread within the country built in, so it can carry an order and has no reason to match the distribution of respondents' answers. This is the representative-agent problem \citep{kirman1992}, and it gives a reason to expect that order and distribution can come apart. Because country-level agreement says nothing about individuals \citep{robinson1950}, we chose the country as the unit and keep our claims there.

Modelling the group directly is not new on its own. SubPOP \citep{suh2025} and Group Preference Optimization \citep{zhao2023gpo} predict group-level responses without simulating individuals, but they train on human response distributions and name the group. What is specific here is the combination of three features, and each is a choice a researcher can inspect. First, the inputs are declared aggregate moments. The researcher states which six GPS moments enter, and a sign rule turns them into labels, so nothing about the conditioning hides in a country name or in human answers. Second, the training data are synthetic preference pairs. Each country's sign profile labels one shared bank of generated response pairs. Direct Preference Optimization \citep{rafailov2023}, which we read as a regularized Bradley--Terry pairwise choice model, then fits one adapter per profile. No human survey answer enters training. Therefore, a human survey is a criterion the method has never seen, which is not true of models fitted to response distributions. Third, the group is unnamed at evaluation. Items are scored with prompts that never mention a country, so whatever the adapters carry comes from the conditioning. A country-token channel is closed by construction.

We evaluate the construction along a chain of three links, treating each as a separate measurement question in the way \citet{wallach2025} propose for generative systems. The anchor is the GPS, the instrument is the set of adapters, and the respondents are the people who answered Wave~7 of the World Values Survey (WVS). Prompted models have been compared with the GPS across countries and found to be weak proxies for it \citep{chen2026}. We built the construction on a purposive development panel of \dc{16}{dev_countries} countries. We then evaluate it on all \dc{42}{countries} countries covered by both the GPS and Wave~7, with a second set of adapters for all \dc{64}{cells} sign profiles on a newly generated scenario bank. At the first link we ask whether training installs the synthetic labels, whether the order of countries survives the new bank, and whether a prompt does as well. A factorial over the \dc{64}{cells} sign profiles asks whether the six declared dimensions behave as separate inputs. At the second link we ran a preregistered screen of \dc{223}{pool_items} Wave~7 items that no adapter was built for. A probe that reads only the adapters' scores assigns each item a dimension and a direction, and a fixed rule selects the strongest assignments before any human value is read. That criterion is retrospective, because the survey existed before the probe. Finally, descriptive analyses added after the registered results had been read set the adapters beside the respondents' composites and response distributions.

What we found on the first link is an order of countries that the declared signs carry. Training installs each country's synthetic labels on scenarios the adapters never saw. On trust, the adapters keep the GPS order of countries when the scenario bank is replaced, whereas a prompt that names the country does not. A prompt that states the same six signs keeps the trust order as well, and the contrast between the adapters and that prompt is unresolved. That is, the signs carry the order, and the evidence gives the trained weights no separate credit. The order is specific to trust. Each declared sign moves its own surface in the intended direction and also moves the others, often by more. Trust is the one surface where its own sign dominates.

Checked against people, the instrument answers unlike them. The retrospective screen gives a weak signal. The registered null is rejected, yet the signal is a small fraction of what two human survey waves achieve on the same items, and its country-bootstrap interval reaches zero. Registered contrasts do not distinguish the adapters from careful readers of the questionnaire on the few items they share, and no dimension passes both the second-bank check and the screen. Descriptively, the adapters' order of countries does not detectably agree with the respondents' composite on any surface. Notably, the GPS's own agreement with those respondents is loose. On response distributions the ranking reverses. The adapters sit farther from the survey's response distributions than the base model does, whereas a country-named prompt sits closer and has no trust order. The separation of order from distributional fidelity is known \citep{dasilva2026,wangc2026}. What this case adds is one method, on one anchor, that answers the two questions in opposite directions. Put plainly, a model asked only about the group can order countries the way a cross-national survey does while answering unlike the people in them.

Our contributions are two.
\begin{enumerate}
\item \textbf{An inspectable group-level conditioning method.} Declared aggregate moments, synthetic preference pairs and an unnamed group at evaluation define a construction whose information budget is on the record. A reader can see what the method was given and check where its readings depart from people.
\item \textbf{A reusable evaluation protocol.} We separate the three links and state what a favorable or unfavorable result at each permits. The protocol includes a screen of new items whose probe is fixed before any human value is read, and a permutation null that absorbs the probe's selection rate. It adds a human ceiling from two survey waves, item-matched comparators that read only the questionnaire, and a factorial that tests whether declared dimensions stay separate.
\end{enumerate}

What the paper shows is bounded by its design: one anchor, one instrument family and one base model. We do not show that trained weights improve on a prompt that states the same signs. Nor do we show that the adapters agree with respondents beyond the weak retrospective signal, or that screening would change which items a researcher fields. Pre-piloting is a prospective use that remains untested. The instrument carries no information about magnitudes or about variation among people within a country, and its readings concern differences among countries. Within these bounds a survey designer can consult it beside human pretesting as a complement, never in its place.

Section~\ref{sec:background} places the method among current uses of synthetic respondents and in the measurement literature on cross-national comparability. Sections~\ref{sec:instrument} and~\ref{sec:design} describe the construction and the design with its registration, Section~\ref{sec:results} reports the results, and Section~\ref{sec:discussion} discusses what they permit.
